\documentclass[10pt]{article}
\usepackage[T1]{fontenc}
\usepackage{lmodern,amsmath,amssymb,amsthm,microtype,booktabs}
\usepackage[margin=0.82in]{geometry}
\usepackage[colorlinks=true,urlcolor=blue,citecolor=blue]{hyperref}
\newtheorem{theorem}{Theorem}
\newtheorem{lemma}{Lemma}
\newtheorem{corollary}{Corollary}
\newcommand{\Z}{\mathbb Z}
\newcommand{\Q}{\mathbb Q}
\newcommand{\vthree}{v_3}
\newcommand{\Stir}[2]{\left\{\begin{smallmatrix}#1\\#2\end{smallmatrix}\right\}}
\title{Powerful Values and 3-Adic Obstructions\\in Cyclic Decimal Concatenations}
\author{Qichao Wang\\{\small Hebei University of Technology}\\[0.3em]
Daoyu Dong\\{\small University of Electronic Science and Technology of China}}
\date{October 26, 2026}
\hypersetup{pdftitle={Powerful Values and 3-Adic Obstructions in Cyclic Decimal Concatenations},pdfauthor={Qichao Wang and Daoyu Dong}}
\begin{document}
\maketitle
\begin{abstract}
We prove a quantitative density-zero bound for powerful values among cyclic
concatenations of initial positive integers, uniformly in the rotation and
in every fixed integer base. At the decimal boundaries $m=10^d-1$, we determine
the exact 3-adic valuation and the normalized residue modulo 9 of every
rotation. These formulas exclude further infinite families of perfect powers.
\end{abstract}

\section{Statements}
In base $g\ge2$, let $W_m$ concatenate $1,2,\ldots,m$, and let $D(m)$
be its digit length, with $W_0=D(0)=0$. For $0\le k<m$, moving the first
$k$ whole integer blocks to the end gives
\begin{equation}\label{eq:rotation}
 R_{m,k}=g^{D(k)}W_m-(g^{D(m)}-1)W_k.
\end{equation}
In base ten these numbers, in increasing order, form OEIS A001292
\cite{oeis}. Kashihara conjectured that none except 1 is a perfect power.
Rip\`a \cite{ripa} excludes $m\equiv2,3,5,6\pmod9$ by a digit-sum argument.
The full conjecture remains open. Here we prove two further restrictions.

Recall that a positive integer is \emph{powerful} if every prime divisor
has exponent at least two. Every nontrivial perfect power is powerful.
Logarithms in the following theorem and in Section~\ref{sec:sieve} are natural.

\begin{theorem}[Quantitative sparsity]\label{thm:sparsity}
For every fixed integer base $g\ge2$,
\begin{equation}\label{eq:sparse}
 \#\{(m,k):2\le m\le X,\ 0\le k<m,\ R_{m,k}\text{ powerful}\}
 \ll_g \frac{X^2(\log\log X)^2}{\log X}.
\end{equation}
For A001292, the corresponding count among the first $N$ terms is
$O(N(\log\log N)^2/\log N)$.
\end{theorem}

For the next result the base is ten and $v_3$ denotes the 3-adic valuation.
\begin{theorem}[Complete digit boundaries]\label{thm:boundary}
Let $d\ge1$, $m=10^d-1$, $h=v_3(d)$, and $u=d/3^h$. Every rotation satisfies
\begin{equation}\label{eq:boundary}
 v_3(R_{m,k})=h+2,\qquad
 \frac{R_{m,k}}{3^{h+2}}\equiv u(2+6d)\pmod9.
\end{equation}
\end{theorem}

\begin{corollary}\label{cor:exclude}
Put $v=h+2$ and let $\mathcal P(v)$ be its set of prime divisors.
No rotation at $m=10^d-1$ is a perfect power in any of the following cases:
\[
\begin{array}{c|c}
 \mathcal P(v)&\text{condition on }u\\ \hline
 \{2\}&u\equiv1\pmod3\\
 \{3\}&u\equiv1,2,7,8\pmod9\\
 \{2,3\}&u\equiv1,7\pmod9.
\end{array}
\]
This includes $d\equiv3,6,9,21,24\pmod{27}$ and $d=81$.
\end{corollary}
\begin{proof}
A prime exponent of a perfect power must divide $v$. A square requires
$v$ even and its normalized unit congruent to 1 modulo 3; by
\eqref{eq:boundary} this requires $u\equiv2\pmod3$.
A cube requires $3\mid v$, hence $h\ge1$ and $3\mid d$.
Its normalized unit must be $\pm1\pmod9$, so $2u\equiv\pm1\pmod9$,
or $u\equiv4,5\pmod9$. The table excludes every available prime exponent.
\end{proof}

\section{Exact 3-adic values at complete boundaries}
Write $V_d=W_{10^d-1}$ and $S(d)=D(10^d-1)$. We will prove the stronger
intermediate congruence
\begin{equation}\label{eq:ratio}
 \frac{V_d}{9d}\in\Z_3,\qquad
 \frac{V_d}{9d}\equiv2+6d\pmod9.
\end{equation}
Quotients by $d$ in this section are taken in $\Q_3$.

\begin{lemma}[Averaging power series]\label{lem:average}
For integers $d\ge1$ and $n\ge1$,
\begin{equation}\label{eq:average}
 v_3\!\left(\frac1d\sum_{j=1}^d j^n\right)
 \ge-\lfloor\log_3(n+1)\rfloor.
\end{equation}
Suppose $f(X)=\sum_{n\ge0}c_nX^n$ has $v_3(c_n)\ge n$,
$c_0\equiv5\pmod9$, $c_1\equiv3\pmod9$, and $v_3(c_2)\ge3$.
Then
\begin{equation}\label{eq:faverage}
 \frac1d\sum_{j=1}^d f(j)\equiv2+6d\pmod9.
\end{equation}
\end{lemma}
\begin{proof}
Using Stirling numbers of the second kind gives the exact identity
\[
 \frac1d\sum_{j=1}^d j^n
 =\sum_{r=0}^n\Stir{n+1}{r+1}
       \frac{r!}{r+1}\binom{d-1}{r}.
\]
Each denominator loses at most $\lfloor\log_3(n+1)\rfloor$ factors of 3,
proving \eqref{eq:average}. For $n\ge3$,
$n-\lfloor\log_3(n+1)\rfloor\ge2$, so all these terms vanish modulo 9
after averaging. The same holds for $n=2$ by its stronger coefficient bound.
Only $c_0+c_1(d+1)/2$ remains, congruent to $5+3(d+1)/2=2+6d\pmod9$.
The coefficient bounds also justify convergence and termwise averaging.
\end{proof}

Define the 3-adic logarithm explicitly by
\[
 \ell=\sum_{r\ge1}\frac{(-1)^{r+1}9^r}{r},\qquad v_3(\ell)=2.
\]
In the following power series put
\[
 E=e^{\ell X},\quad a=E/10,\quad q=9a,\quad b=E-1,\quad
 S(X)=\frac{(9X-1)E+1}{9},
\]
and define
\begin{equation}\label{eq:U}
 \frac{U(X)}9=a^2+
 \frac{ab+1}{9}\sum_{r\ge0}\binom q{r+2}b^r,
 \qquad f(X)=\frac{U(X)}9e^{-\ell S(X)}.
\end{equation}
At every positive integer $j$, $S(j)$ is the digit length of $V_j$,
and $U(j)$ concatenates all $j$-digit integers. Indeed, with
$a=10^{j-1}$, $b=10^j-1$, and $q=9a$, the usual arithmetic-geometric sum gives
\[
 b^2U(j)=(ab+1)(1+b)^q-((1+b)b+1).
\]
Expanding $(1+b)^q$ cancels its constant and linear terms and gives
\eqref{eq:U}, since $a+q=10^j$ and $aq=9a^2$.

\begin{lemma}\label{lem:coefficients}
The coefficients of $f$ satisfy all hypotheses of Lemma~\ref{lem:average}.
\end{lemma}
\begin{proof}
Let $\mathcal A$ consist of series whose degree-$n$ coefficient has
valuation at least $n$. It is closed under addition and multiplication.
In the $r$th summand of $U/9$, the factor $b^r$ makes all degrees below $r$
zero. The factor $q$ in $\binom q{r+2}$ cancels the displayed denominator 9.
The remaining expression is a polynomial in $E$ with 3-adic denominators
bounded by $(r+2)!$. Its degree-$n$ coefficient, for $r\le n$, therefore has
valuation at least
\[
 2n-v_3(n!)-v_3((r+2)!)\ge n.
\]
Here $v_3(n!)\le(n-1)/2$ for $n\ge1$, and
$v_3((r+2)!)\le(n+1)/2$; the degree-zero case is immediate.
The term $a^2$ satisfies the same bound. Thus $U/9\in\mathcal A$.

The constant term of $S$ is zero. Its linear coefficient is $1-\ell/9$,
of valuation at least 2, and for $n\ge2$ its coefficient is
$\ell^{n-1}(n-\ell/9)/n!$, of valuation at least $n$.
Hence $S\in\mathcal A$. Expanding $e^{-\ell S}$ preserves these bounds:
its $r$th exponential term gains $2r-v_3(r!)\ge0$ factors of 3.
Consequently $f\in\mathcal A$.

Expansion through degree two in \eqref{eq:U} gives
\begin{align*}
 c_0&=\frac1{200},\\
 c_1&=\frac{169\ell}{3000}+\frac{\ell^2}{1800},\\
 c_2&=\frac{671\ell^2}{80000}
        +\frac{353\ell^3}{54000}+\frac{\ell^4}{32400}.
\end{align*}
Since $\ell/9\equiv1\pmod9$, these imply
$c_0\equiv5\pmod9$, $c_1\equiv3\pmod9$, and $v_3(c_2)\ge3$.
\end{proof}

\begin{proof}[Proof of Theorem~\ref{thm:boundary}]
The concatenation of the complete digit blocks yields
\[
 V_d=10^{S(d)}\sum_{j=1}^d U(j)10^{-S(j)}
     =9\,10^{S(d)}\sum_{j=1}^d f(j).
\]
Lemma~\ref{lem:average} and $10^{S(d)}\equiv1\pmod9$ prove
\eqref{eq:ratio}. Its right side is a 3-adic unit, so $v_3(V_d)=h+2$.

To transfer the normalized unit to every rotation, note that
\begin{equation}\label{eq:S}
 S(d)=\frac{9d(d+1)}2
 \left(1+\sum_{r=1}^{d-1}\frac{2\cdot9^r}{r+2}\binom{d-1}{r}\right).
\end{equation}
This follows by expanding $S(d)=9\sum_{j=1}^d j10^{j-1}$ in powers of 9.
Each term after 1 in the bracket is divisible by 3; hence
$v_3(S(d))=2+h+v_3(d+1)$. Lifting the exponent gives
\[
 v_3(10^{S(d)}-1)=4+h+v_3(d+1).
\]
In \eqref{eq:rotation}, $10^{D(k)}\equiv1\pmod9$ and the second term
is divisible by $3^{h+4}$. Thus $R_{m,k}\equiv V_d\pmod{3^{h+4}}$.
Combining this with \eqref{eq:ratio} proves both assertions in
\eqref{eq:boundary}.
\end{proof}

\section{A uniform sieve for powerful values}\label{sec:sieve}
We now prove Theorem~\ref{thm:sparsity}. In a fixed $d$-digit interval set
\[
 B=g^d,\quad L=g^{d-1},\quad c=B-1,\quad
 A=c^2W_{L-1}+Lc+1.
\]
The recurrence $W_{m+1}=BW_m+(m+1)$ gives, for $L-1\le m\le B-1$,
\[
 c^2W_m=A B^{m-L+1}-((m+1)c+1),\qquad
 D(m)=D(L-1)+d(m-L+1).
\]
For fixed $k$, writing $u=g^{D(k)}$, substitution into
\eqref{eq:rotation} gives
\begin{equation}\label{eq:affine}
 R_{m,k}=\alpha B^m+\beta m+\gamma,
\quad
 \begin{cases}
 \alpha=(uA/c^2-g^{D(L-1)}W_k)B^{1-L},\\
 \beta=-u/c,\quad\gamma=W_k-uB/c^2.
 \end{cases}
\end{equation}
At every prime $p\nmid g(B-1)$ the coefficients are defined modulo $p^2$
and $\beta$ is a unit. The right side provides a periodic modular extension;
only its values inside the digit interval are identified with the concatenation.

\begin{lemma}[Valuation-one densities]\label{lem:densities}
Let $p\nmid B\beta$ and $F(n)=\alpha B^n+\beta n+\gamma\pmod{p^2}$.
The event $p\mid F(n)$, $p^2\nmid F(n)$ has period dividing
$T_p=p^2(p-1)$ and density $\rho_p=(p-1)/p^2$.
For any finite set $P$ of eligible primes the events have joint density
$\prod_{p\in P}\rho_p$ in a period dividing $\operatorname{lcm}_{p\in P}T_p$.
\end{lemma}
\begin{proof}
Fix $n\pmod{p-1}$. The exponential term is constant modulo $p$, and the
linear coefficient is nonzero, so exactly one class $n\pmod p$ is a root.
Its $p$ lifts differ by $p(p-1)t$. Euler's theorem fixes the exponential
term modulo $p^2$ on these lifts, while the linear term takes $p$ distinct
values. Precisely one lift is zero modulo $p^2$, leaving $p-1$ of valuation one.

For the joint assertion induct on the largest prime $q$.
The smaller-prime events have a period $T$ with no prime divisor as large
as $q$. Thus $q\nmid H=\operatorname{lcm}(T,q-1)$.
In each fixed class modulo $H$, the smaller-prime events are fixed and
the Chinese remainder theorem lets $n\pmod{q^2}$ vary freely. The local
count gives the same fraction $\rho_q$ in every class. Averaging multiplies
the previous joint density by $\rho_q$. This allows overlapping periods,
including a smaller prime dividing a larger prime minus one.
\end{proof}

For an interval $I$ of $H$ integers and a subset $P'$, Lemma~\ref{lem:densities}
gives its simultaneous event count as
$H\prod_{p\in P'}\rho_p+O(\prod_{p\in P'}p^3)$, since the interval error
is at most the joint period. Take eligible primes at most $z$, put
$\lambda=\sum\rho_p$, and choose an even integer $K\ge0$.
Powerful values avoid every valuation-one event. Finite Bonferroni gives
\begin{equation}\label{eq:sieve}
 \#\{n\in I:F(n)\text{ powerful}\}
 \le H\left(e^{-\lambda}+\frac{\lambda^{K+1}}{(K+1)!}\right)
       +(K+1)z^{4K}.
\end{equation}
Indeed, if $e_j$ are the elementary symmetric sums of the $\rho_p$,
the even truncated main term is at most
$\prod(1-\rho_p)+e_{K+1}\le e^{-\lambda}+\lambda^{K+1}/(K+1)!$.
There are at most $z^j$ subsets of size $j$, each costing at most $z^{3j}$,
which proves the displayed error. These estimates are uniform in the
coefficients of \eqref{eq:affine} and in $k$.

Choose
\[
 \eta=\frac{\log g}{100},\qquad
 z=\exp\!\left(\frac{\eta d}{\log d}\right),\qquad
 K=2\lceil4\log d\rceil,
\]
and use primes $d^2<p\le z$ with $p\nmid g(g^d-1)$.
For large $d$, $z>d^2$. Omitted primes exceeding $d^2$ and dividing $g^d-1$
number at most $\log(g^d-1)/(2\log d)$, with reciprocal sum
$O_g(1/(d\log d))$. The fixed prime divisors of $g$ eventually lie below $d^2$.
Mertens' estimate \cite{mertens} therefore yields
\[
 \lambda=\log\log z-\log\log(d^2)+O_g(1)
        =\log d-2\log\log d+O_g(1).
\]
In particular $0<\lambda\le\log d$ eventually, so
$e^{-\lambda}\ll_g(\log d)^2/d$. Also
\[
 \frac{\lambda^{K+1}}{(K+1)!}\le(e/8)^{K+1}\le d^{-8},\qquad
 (K+1)z^{4K}<g^{d/2}
\]
for large $d$; the latter follows from $K\le10\log d$ and
$\log((K+1)z^{4K})\le\log(K+1)+40\eta d<d\log g/2$.
Thus for every fixed $k$ and every interval inside the digit block with $m>k$,
\begin{equation}\label{eq:uniform}
 \#\{m\in I:R_{m,k}\text{ powerful}\}
 \ll_g |I|\frac{(\log d)^2}{d}+g^{d/2}.
\end{equation}

Let $D$ be the digit length of $\lfloor X\rfloor$.
For each $d\le D$ and $k$, apply \eqref{eq:uniform} on
$\max(g^{d-1},k+1,2)\le m\le\min(X,g^d-1)$.
There are at most $g^d$ choices of $k$, so the total interval error is
$\sum_{d\le D}g^{3d/2}=O_g(X^{3/2})$.
Blocks with $d\le D/2$ contain only $O_g(X)$ pairs and may be bounded trivially.
On the other blocks the main factor is $O((\log D)^2/D)$, and there are
$O(X^2)$ pairs. This proves \eqref{eq:sparse}.

Finally, each $m$ supplies $m$ distinct rotations. If $m$ has $d\ge2$ digits,
the cyclic digit word has a unique occurrence of 1 followed by $d-1$ zeros:
it is the integer $g^{d-1}$. Such a zero run cannot cross an integer boundary
or the cyclic join, whose next digit is 1. The unique marker prevents a
nonidentity digit rotation from fixing the word. For $d=1$ all digits are distinct.
Digit lengths increase strictly with $m$, so the sequence through block $m=X$
has $X(X+1)/2$ entries. Bounding a partial final block by the whole block
and taking $X\asymp\sqrt N$ proves the sequence-index assertion.
\qed

\begin{thebibliography}{9}
\bibitem{oeis} OEIS Foundation, Sequence A001292,
\url{https://oeis.org/A001292}, accessed October 26, 2026.
\bibitem{ripa} M. Rip\`a, \emph{On some open problems concerning perfect powers},
arXiv:2205.10163, version 4 (2024), \url{https://arxiv.org/abs/2205.10163}.
\bibitem{mertens} F. Mertens, Ein Beitrag zur analytischen Zahlentheorie,
\emph{J. reine angew. Math.} \textbf{78} (1874), 46--62,
\url{https://doi.org/10.1515/crll.1874.78.46}.
\end{thebibliography}
\end{document}
